\documentclass{situationist-preprint}

\title{Against Philosopher-Kings}
\msid{POL-15-a}
\msfamily{POL-15}
\msdiscipline{Political Philosophy}
\mstype{Principal manuscript}
\mskeywords{philosopher-kings; expertise; political authority; epistocracy; Plato; expert/boss fallacy}

\begin{document}
\maketitle

\begin{abstract}
Should philosophers rule? Plato argued that cities will have no rest from evils until philosophers become kings, because only those who know the good can govern well. We argue that philosophers should not rule, because philosophical expertise does not confer political authority. The argument has two steps. First, even if some people know better than others what ought to be done, it does not follow that they have the right to command others to do it: authority is a normative relation between persons, and knowledge is not. This is the gap that Estlund calls the expert/boss fallacy. Second, the only route across the gap, that it would be reasonable for those subject to the expert to defer, is closed in the political case, because the identification of philosophical expertise on questions of justice is subject to reasonable disagreement, and because deference on such questions cannot be confined to the questions on which the expert is reliable. Philosophers may advise, argue and persuade; they have no claim to rule.
\end{abstract}

\section{The question}

In the \emph{Republic}, Socrates says that unless philosophers rule as kings, or those now called kings genuinely philosophize, there will be no end to the troubles of cities \citep[473c--d]{PlatoRep}. The argument rests on an analogy with crafts: as the ship should be steered by the navigator and the patient treated by the physician, the city should be governed by those who know the good \citep[488a--489c]{PlatoRep}. \citet{Popper1945} read the proposal as the charter of closed and authoritarian society. The contemporary descendant of the argument is epistocracy, the view that political power should be distributed according to competence \citep{Brennan2016}. Our question is narrower: does philosophical expertise, supposing there is such a thing, confer the authority to rule?

\section{Expertise and authority}

\begin{definition}\label{d:expertise}
A person has \emph{expertise} in a domain if her judgments in that domain are more reliable than those of most others.
\end{definition}

\begin{definition}\label{d:authority}
A person has \emph{political authority} over others if her directives give them reasons to comply that they would not otherwise have, and she has a right, enforceable by coercion, that they comply.
\end{definition}

\begin{proposition}[The gap]\label{p:gap}
Expertise in a domain does not by itself confer authority over others in that domain.
\end{proposition}

\begin{proof}
Expertise is a property of a person's judgments; authority is a relation between persons. An expert's judgment that some act ought to be done is, if she is reliable, evidence that it ought to be done, and may give others reason to believe it. It does not give them a reason to do it simply because she directs it, nor does it give her a right that they do it. A physician who knows that a patient should take a drug has no right to compel him to take it. That one knows better, as \citet{Estlund2008} puts it, does not make one the boss.
\end{proof}

The gap can be crossed only by a further premise. The most promising is Raz's \emph{normal justification thesis}: a person has authority over another if the other would better conform to the reasons that apply to him by following her directives than by acting on his own judgment \citep{Raz1986}. On this thesis expertise matters, because the expert's directives are more likely to track the applicable reasons.

\begin{proposition}\label{p:closed}
The normal justification thesis does not confer political authority on philosophers.
\end{proposition}

\begin{proof}
Three conditions would have to hold. (i) The philosopher's directives would have to be more reliable than citizens' own judgments on the questions of political rule. (ii) Citizens would have to be able to identify which philosophers are the reliable ones. (iii) The authority would have to be limited to the questions on which the philosopher is reliable. Condition (i) is doubtful: political rule requires knowledge of particular facts, of institutions and of consequences, in which philosophical training confers no special reliability, and on the normative questions themselves philosophers disagree as much as anyone. Condition (ii) fails in the political case: the identification of experts on justice requires a view about justice, and reasonable citizens hold different views, so any designation of philosopher-rulers is one that some reasonable citizens can reject \citep{Estlund2008}. Condition (iii) fails because political rule is general: a ruler decides what the questions are and which considerations bear on them, and the authority to rule cannot be restricted in advance to questions on which the ruler has shown reliability. Since Raz's thesis grounds authority only within the domain in which following the directives improves conformity with reason, it grounds no general authority to rule.
\end{proof}

\begin{theorem}\label{t:main}
Philosophers should not rule: philosophical expertise does not confer political authority.
\end{theorem}

\begin{proof}
By Proposition~\ref{p:gap}, expertise alone does not confer authority. By Proposition~\ref{p:closed}, the strongest available principle for converting expertise into authority does not apply to philosophers in the political case. Without a ground of authority, a claim to rule is a claim to coerce without right.
\end{proof}

\section{Objections}

\begin{objection}[Plato's analogy]
We accept the authority of the navigator on a ship and of the surgeon in the operating theater. The city is a ship.
\end{objection}

\begin{reply}
The navigator's and the surgeon's authority rests on the passengers' and patient's prior choice of the end, safe arrival or health, and on their consent to the expert's management of means. The city has no end agreed in advance; what its end should be is the question in dispute. Where the end is contested, the analogy assigns the expert authority over precisely what the analogy assumes settled.
\end{reply}

\begin{objection}[Benevolent rule]
If philosopher-rulers would in fact produce better outcomes, the citizens would be better off, and that is reason enough.
\end{objection}

\begin{reply}
Better outcomes can justify adopting a procedure only if the procedure is otherwise permissible. A procedure that ranks citizens as fit and unfit to decide their common affairs denies them equal standing, and good outcomes alone do not justify a procedure that does so (POL-04).
\end{reply}

\begin{objection}[Advice without authority is idle]
If philosophers have no authority, their expertise has no political role.
\end{objection}

\begin{reply}
It has the role it has in any free society: argument in public, advice to officials and citizens, and the clarification of principles. These give others reasons to believe, which they may weigh for themselves. That is the role compatible with Proposition~\ref{p:gap}.
\end{reply}

\section{Conclusion}

Even if philosophers knew better than others what justice requires, knowledge does not confer the right to rule, and the principle that would convert expertise into authority cannot do so for political rule, where experts cannot be identified without taking sides in the dispute they are meant to settle and where authority cannot be confined to the domain of expertise. Philosophers should not rule.

\begin{thebibliography}{9}\small
\bibitem[Brennan(2016)]{Brennan2016} Brennan, J. (2016). \emph{Against Democracy}. Princeton University Press.
\bibitem[Estlund(2008)]{Estlund2008} Estlund, D. (2008). \emph{Democratic Authority: A Philosophical Framework}. Princeton University Press.
\bibitem[Plato(Rep.)]{PlatoRep} Plato. \emph{Republic}.
\bibitem[Popper(1945)]{Popper1945} Popper, K. R. (1945). \emph{The Open Society and Its Enemies}, vol.~1: \emph{The Spell of Plato}. Routledge.
\bibitem[Raz(1986)]{Raz1986} Raz, J. (1986). \emph{The Morality of Freedom}. Oxford University Press.
\end{thebibliography}

\end{document}
